\documentclass{article}

% --- Paquetes solicitados ---
\usepackage[T1]{fontenc}
\usepackage[utf8]{inputenc}
\usepackage[spanish,es-tabla]{babel}
\usepackage[margin=2.4cm]{geometry}
\usepackage{amsmath}
\usepackage{amssymb}
\usepackage{booktabs}
\usepackage{array}
\usepackage{enumitem}
\usepackage{microtype}

% --- Información del documento ---
\title{Documento de Prueba de Paquetes LaTeX}
\author{Prueba de Configuración}
\date{\today}

\begin{document}

\maketitle

\section{Introducción}
Este documento sirve para verificar que los paquetes solicitados carguen correctamente y funcionen sin conflictos. A continuación, se muestra una demostración práctica de las capacidades de cada uno.

\section{Demostración de Paquetes}

\subsection{Texto y Tipografía (\texttt{babel}, \texttt{fontenc}, \texttt{inputenc}, \texttt{microtype})}
Gracias a \texttt{babel} con la opción \texttt{es-tabla}, las construcciones automáticas del lenguaje se adaptan al español (por ejemplo, la fecha de hoy o el nombre de las secciones). Además, podemos escribir eñes (ñ), acentos (á, é, í, ó, ú) y signos de apertura ($\text{¿, ¡}$) directamente en el código fuente. El paquete \texttt{microtype} mejora sutilmente el espaciado y la justificación de este párrafo para evitar saltos de línea extraños.

\subsection{Matemáticas Avanzadas (\texttt{amsmath}, \texttt{amssymb})}
El paquete \texttt{amsmath} permite estructurar ecuaciones complejas, mientras que \texttt{amssymb} proporciona símbolos matemáticos adicionales como los conjuntos numéricos $\mathbb{R}$, $\mathbb{Z}$ o $\mathbb{N}$.

Una ecuación alineada de prueba:
\begin{align}
    f(x) &= \int_{a}^{b} g(x) \, dx \\
    \sum_{i=1}^{n} i &= \frac{n(n+1)}{2}
\end{align}

\subsection{Tablas Estilizadas (\texttt{booktabs}, \texttt{array})}
El paquete \texttt{booktabs} elimina las líneas verticales innecesarias y genera tablas de aspecto profesional usando comandos como \texttt{\textbackslash toprule}, \texttt{\textbackslash midrule} y \texttt{\textbackslash bottomrule}. El paquete \texttt{array} mejora el formato de las celdas.

\begin{table}[h]
    \centering
    \caption{Ejemplo de tabla profesional (comprobar que dice ``Tabla'' y no ``Cuadro'').}
    \vspace{0.3cm}
    \begin{tabular}{p{3cm} m{4cm} c}
        \toprule
        \textbf{Paquete} & \textbf{Función Principal} & \textbf{Estado} \\
        \midrule
        \texttt{geometry} & Control de márgenes (2.4 cm) & Operativo \\
        \texttt{booktabs} & Líneas de calidad editorial & Operativo \\
        \texttt{enumitem} & Personalización de listas & Operativo \\
        \bottomrule
    \end{tabular}
    \label{tab:ejemplo}
\end{table}

\subsection{Listas Personalizadas (\texttt{enumitem})}
Con \texttt{enumitem} podemos modificar el espaciado y el tipo de viñeta de las listas fácilmente mediante opciones en corchetes.

\begin{itemize}[label=$\rightarrow$, noitemsep]
    \item Primer elemento con flecha matemática.
    \item Segundo elemento con espacio reducido.
    \item Tercer elemento de prueba.
\end{itemize}

\end{document}
